% Part: lambda-calculus
% Chapter: introduction
% Section: computable-lambda

\documentclass[../../../include/open-logic-section]{subfiles}

\begin{document}

\olfileid{lam}{int}{lrp}
\olsection{الدوال القابلة للحساب \usetoken{S}{lambda definable}}

\begin{thm}
\ollabel{thm:computable-lambda}
كل دالة جزئية تقبل الحساب هي !!{lambda definable}.
\end{thm}

\begin{proof}
المطلوب أن تكون كل دالة جزئية قابلة للحساب~${\OLId{latin-ordinary}{f}}$
!!{lambda defined} بحد لامبدا~${\OLId{latin-looped}{F}}$.
وترد مبرهنة كليني للصورة السوية هذا المطلوب إلى أمرين:
أن كل دالة عودية بدائية !!{lambda defined} بحد لامبدا،
وأن الدوال !!{lambda definable} مغلقة تحت ما يلزم من
التركيب والبحث غير المحدود.
ولإثبات أن كل دالة عودية بدائية !!{lambda defined} بحد لامبدا،
يكفينا إثبات أن الدوال الابتدائية !!{lambda definable}،
وأن الدوال الجزئية !!{lambda definable} مغلقة تحت التركيب
والعودية البدائية والبحث غير المحدود.
\end{proof}

ونقرب بقية البرهان إلى القراءة باستعمال ترميز مألوف؛
فنكتب مثلًا~${\OLId{latin-looped}{M}}({\OLId{latin-ordinary}{x}}, {\OLId{latin-ordinary}{y}}, {\OLId{latin-ordinary}{z}})$ عوض~$Mxyz$.
لكن لا يحملنّ هذا الترميز على أن لحدود حساب لامبدا غير
المنمّط أعداد وسائط ثابتة: يجوز في الحد نفسه~${\OLId{latin-looped}{M}}$
أن نكتب~${\OLId{latin-looped}{M}}({\OLId{latin-ordinary}{x}},{\OLId{latin-ordinary}{y}})$ وأن نكتب~${\OLId{latin-looped}{M}}({\OLId{latin-ordinary}{x}},{\OLId{latin-ordinary}{y}},{\OLId{latin-ordinary}{z}},{\OLId{latin-ordinary}{w}})$.
وإنما يفيد الترميز أننا نعامل~${\OLId{latin-looped}{M}}$ معاملة دالة في ثلاثة
متغيرات، فيظهر به قصد التعريف.
وبمثل ذلك نقول: «عرّف~${\OLId{latin-looped}{M}}$ بالمعادلة~${\OLId{latin-looped}{M}}({\OLId{latin-ordinary}{x}},{\OLId{latin-ordinary}{y}},{\OLId{latin-ordinary}{z}}) = \dots$»،
عوض قولنا: «عرّف~${\OLId{latin-looped}{M}}$ بالحد~${\OLId{latin-looped}{M}} = \lambd[{\OLId{latin-ordinary}{x}}][\lambd[{\OLId{latin-ordinary}{y}}][\lambd[{\OLId{latin-ordinary}{z}}][\dots]]]$».

\end{document}
